\documentclass[crop=false]{standalone}
% --- ПРЕАМБУЛА ДЛЯ ПОДДЕРЖКИ РУССКОГО ЯЗЫКА И МАТЕМАТИКИ ---
\usepackage[T2A]{fontenc}
\usepackage[utf8]{inputenc}
\usepackage[russian]{babel}
\usepackage{amsmath}
\usepackage{geometry}
\usepackage{amssymb}
\usepackage{graphicx}
\usepackage{float}
\usepackage{import}
\geometry{a4paper, top=2cm, bottom=2cm, left=2cm, right=2cm}
\newcommand{\sidecomment}[1]{\marginpar{\raggedright\scriptsize\itshape #1}}
\newcommand{\iu}{\mathrm{i}}

\newcommand{\Def}{%
    \text{Def}%
    \hspace{-0.85em}% Сдвиг назад ровно на середину (подберите под шрифт)
    \raisebox{0.1ex}{/}% Черта
    \hspace{1em}% Возврат к концу слова + небольшой отступ
}

\renewcommand{\Re}{\operatorname{Re}}
\renewcommand{\Im}{\operatorname{Im}}

\ifstandalone
    \newcommand{\lecturebreak}{} % В отдельном файле лекции разрыв не нужен
\else
    \newcommand{\lecturebreak}{\clearpage} % В полной книге - начинаем с новой страницы
\fi

\newcounter{lecturenumber} % Создаем счетчик для нумерации лекций
\newcommand{\lecture}[2]{
    \lecturebreak
    \refstepcounter{lecturenumber} % Увеличиваем счетчик на 1
    \par\vspace{2em}\noindent % Отступ сверху и убираем абзацный отступ
    {\Large\bfseries % Делаем заголовок крупным и жирным
    Лекция \thelecturenumber: #1 % Выводим "Лекция N: Тема"
    \hfill % Растягиваем пробел, чтобы дата была справа
    \large\textit{#2}} % Выводим дату курсивом
    \par\vspace{1em} % Небольшой отступ снизу
    \addcontentsline{toc}{section}{Лекция \thelecturenumber: #1} % Добавляем в оглавление
    \par
}

\newcommand{\TwoCols}[2]{%
  \par\noindent % Начать с новой строки без отступа
  \begin{minipage}[t]{0.48\textwidth}
    #1
  \end{minipage}% <-- Важный '%' для удаления пробела
  \hfill % Растягивающийся пробел между колонками
  \begin{minipage}[t]{0.48\textwidth}
    #2
  \end{minipage}% <-- Важный '%' для удаления пробела
  \par%
}

\pagestyle{plain} % Включаем нумерацию страниц\
\begin{document}
\lecture{}{15.05}
\[\frac{d}{d t} \frac{\partial T}{\partial \dot q} - \frac{\partial T}{\partial q_i} = Q_i = - \frac{\partial \Pi}{\partial q_i} + Q_i^* \qquad \frac{d}{d t} \frac{\partial L}{\partial \dot q} - \frac{\partial L}{\partial q_i} = Q_i^*\]
\begin{theorem}
    \[T = T_2 + T_1 + T_0\]
    где $T_2, T_1, T_0$ — однордные многочлены соответствующей степени относительно обобщённых скоростей
\end{theorem}
\begin{proof}
    \[
    \begin{aligned}
    T = \frac{1}{2} \sum_{k=1}^{N} m_k v_k^2 &= \frac{1}{2} \sum_{k=1}^{N} m_k \left( \sum_{i=1}^{n} \frac{\partial \bar r_k}{\partial q_i} \dot q_i + \frac{\partial \bar r_k}{\partial t}\right)^2 \\
    &=\frac{1}{2} \sum_{k=1}^{N} m_k \left( \sum_{i, j=1}^{n} \frac{\partial \bar r_k}{\partial q_i} \frac{\partial \bar r_k}{\partial q_j} \dot q_i \dot q_j + 2 \sum_{i=1}^{n}  \frac{\partial \bar r_k}{\partial q_i} \frac{\partial \bar r_k}{\partial t} \dot q_i + \left(\frac{\partial \bar r_k}{\partial t}\right)^2 \right) \\
    &= \underbracket{\frac{1}{2} \sum_{i,j = 1}^{n} \sum_{k=1}^{N} m_k \frac{\partial \bar r_k}{\partial q_i} \frac{\partial \bar r_k}{\partial q_j} \dot q_i \dot q_j}_{T_2}
    + \underbracket{\sum_{i = 1}^{n} \sum_{k=1}^{N} m_k \frac{\partial \bar r_k}{\partial q_i} \frac{\partial \bar r_k}{\partial t} \dot q_i }_{T_1}
    + \underbracket{\frac{1}{2} \sum_{k=1}^{N} m_k \left(\frac{\partial \bar r_k}{\partial t}\right)^2}_{T_0}
    \end{aligned}
    \]
    где $T_2, T_1, T_0$ — однородные многочлены соответствующих порядков
    \[\begin{array}{l}
        P(\lambda x) = \lambda^k P(x) \\ \\ T = T_2 + T_1 + T_0
    \end{array}
        \qquad
    \begin{dcases}
        T_2 = \frac{1}{2} \sum_{i, j = 1}^{n} a_{i j} \dot q_i \dot q_j, & a_{i j} = \sum_{k=1}^{N} m_k \frac{\partial \bar r_k}{\partial q_i} \frac{\partial \bar r_k}{\partial q_j} \\
        T_1 = \sum_{i=1}^{n} a_i \dot q_i, & a_i = \sum_{k=1}^{N} m_k \frac{\partial \bar r_k}{\partial q_i} \frac{\partial \bar r_k}{\partial t} \\
        T_0 = \frac{1}{2} \sum_{k=1}^{N} m_k \left(\frac{\partial \bar r_k}{\partial t}\right)^2
    \end{dcases}
    \]
\end{proof}

Замечание: Так как потенциальная энергия не зависит от скоростей и ускорений:
\[L = T - \Pi = T_2 + T_1 + \underbracket{T_0 - \Pi}_{\tilde T_0}\]

\noindent \keypoint 
Уравнения лагранжа образуют СЛАУ по обобщённым ускорениям
\begin{proof}   
    \[
    \begin{aligned}
    &\frac{d}{d t} \frac{\partial (T_2 + T_1 + T_0)}{\partial \dot q} - \frac{\partial T}{\partial q_i} = Q_i \qquad &\frac{\partial T_0}{\partial \dot q_i} = 0 \qquad \frac{\partial T_1}{\partial \dot q_i} = a_i \qquad \frac{\partial T_2}{\partial \dot q_i} = \sum_{j=1}^{n} a_{i j } \dot q_j \\
    &\frac{d}{d t} \left( \sum_{j=1}^{n} a_{i j} \dot q_j+ a_i \right) = Q_i + \frac{\partial T}{\partial q_i} \qquad &\frac{d}{d t} \left( \sum_{j=1}^{n} a_{i j} \dot q_j\right) = Q_i - \frac{d a_i}{d t}+ \frac{\partial T}{\partial q_i}
    \end{aligned}\]
    \[\sum_{j=1}^{n} a_{i j} \ddot q_j = - \sum_{j=1}^{n} \dot q_j \dot a_{i j} - \dot a_i + Q_i + \frac{\partial T}{\partial q_i} := G_i(q_1 , \dots , q_n, \dot q_1, \dots, \dot q_n, t)\]
    \[\sum_{j=1}^{n} a_{i j} \ddot q_j = G_i\]
\end{proof}
    \begin{theorem}[Основная теорема Лагранжевого формализма]
        \[\det (a_{i j})_{i, j = 1}^n \neq 0 \qquad \forall q, t\]
    \end{theorem}
    \begin{proof}
    \[a_{i j} = m_1 \left( \frac{\partial x_1}{\partial q_i} \frac{\partial x_1}{\partial q_j} + \frac{\partial y_1}{\partial q_i} \frac{\partial y_1}{\partial q_j} + \frac{\partial z_1}{\partial q_i} \frac{\partial z_1}{\partial q_j}\right) + \dots\]
    Теорема Грама(взяли подматрицу  из матрицы Грама), если коэффициенты это скалярные произведения л.н.з векторов, то матрица приведённой системы невырождена
    \[a_{i j} = \bar u_i \bar u_j\]
    \[u_i' = \left(\frac{\partial x_1}{\partial q_i} \quad \frac{\partial y_1}{\partial q_i} \dots \frac{\partial z_N}{\partial q_i}\right)^T\]
    \[
    \underset{3 N \times n}{U}= \begin{pmatrix}
        \frac{\partial x_1}{\partial q_1} & \frac{\partial x_1}{\partial q_2} & \frac{\partial x_1}{\partial q_3} \dots \frac{\partial x_1}{\partial q_n} \\
        \frac{\partial y_1}{\partial q_1} & \frac{\partial y_1}{\partial q_2} & \frac{\partial y_1}{\partial q_3} \dots \frac{\partial y_1}{\partial q_n} \\
        \frac{\partial z_1}{\partial q_1} & \vdots \\
        \frac{\partial x_2}{\partial q_1} & \vdots \\
        \vdots \\
        \frac{\partial z_N}{\partial q_1} & \frac{\partial z_N}{\partial q_2} & \frac{\partial z_N}{\partial q_3} \dots \frac{\partial z_N}{\partial q_n} \\
    \end{pmatrix} \begin{matrix}
        \sqrt m_1 \\ \sqrt m_1 \\ \sqrt m_1 \\ \sqrt m_2  \\ \vdots \\  \sqrt m_N % fixme поправить
    \end{matrix}
    \]
    Столбцы матрица линейно независимы в силу независимости обобщённых координат, то есть и векторы составленные из данных координатных столбцов линейно независимы(домножение на константу это элементарное преобразование)
\end{proof}



\Def Натуральные системы — системы, в которых лагранжиан является разностью кинетической и потенциальной энергии

\Def Ненатуральные системы — не натуральные системы(в таких надо потребовать невырожденность матрицы сосставленной из $a_{i j} = \frac{\partial^2 L}{\partial \dot q_i \partial \dot q_j}$ при допустимых $q_i , q_j$)


\noindent \keypoint Т Эйлера об однородной функции

для $F(\lambda x_1, \lambda x_2, \dots, \lambda x_n) = \lambda^k F(x_1, x_2, \dots, x_n)$
\[\sum_{i=1}^{n} \frac{\partial F}{\partial x_i} x_i = k F\]

\begin{theorem}[Об изменении полной энергии голономной системы]
    \[\frac{d E}{d t} = \sum_{i=1}^{n} Q_i^* \dot q_i + \frac{d}{d t} (T_1 + 2 T_0) + \frac{\partial \Pi}{\partial t} - \frac{\partial T}{\partial t}\]
    Мощность непотенциальных сил, частная производная потенциальной энергии и вклад нестационарных связей

\end{theorem}
\begin{proof}
    \[
        \begin{aligned}
            \frac{d}{d t} T &= \frac{d}{d t}\sum_{i=1}^{n} \frac{\partial T}{\partial \dot q_i} \dot q_i - \sum_{i=1}^{n} \frac{d}{d t}\left(\frac{\partial T}{\partial \dot q_i} \right)\dot q_i + \sum_{i=1}^{n}  \frac{\partial T}{\partial q_i} \dot q_i +\frac{\partial T}{\partial t}  \\
            &= \frac{d}{d t} \sum_{i=1}^{n} \frac{\partial T}{\partial \dot q_i} \dot q_i  - \sum_{i=1}^{n}\left(\frac{d}{d t} \frac{\partial T}{\partial \dot q_i} - \frac{\partial T}{\partial q_i} \right) \dot q_i + \frac{\partial T}{\partial t} = \frac{d}{d t} (2 T_2 + T_1) - \sum_{i=1}^{n} Q_i^* \dot q_i + \sum_{i=1}^{n} \frac{\partial \Pi}{\partial q_i} \dot q_i+ \frac{\partial T}{\partial t} \\
            &= 2 \frac{d T}{d t} - \frac{d}{d t} (T_1 + 2 T_0) - \sum_{i=1}^{n} Q_i^* \dot q_i  + \frac{d \Pi}{d t} - \frac{\partial \Pi}{\partial t} + \frac{\partial T}{\partial t}
        \end{aligned}
        \]
        \[ E = T + \Pi\]
    \end{proof}
Следствие: В консервативных системах:
\[\frac{d E}{d t} = 0\]
\end{document}